\begin{topic}[id=neuroprostheses,
             difficulty={Anspruchsvoll},
             type={Literaturreview + Einordnung},
             prerequisites={Grundlagen Signalverarbeitung; Offenheit für interdisziplinäre Literatur},
             practice={optional},
             owner={Dozententeam},
             updated={2026-04-09},
             tags={ Neuroprothesen, BCI, Stimulation, Signalverarbeitung, Ethik }]{Neuroprothesen}

\TopicDescription{Neuroprothesen koppeln Nervensystem und Technik: von Cochlea-Implantaten bis zu motorischen BCIs. Wir diskutieren Signalwege (invasiv/nichtinvasiv), Stimulation/Decoding, Sicherheit/Ethik und klinische Rahmenbedingungen. Ziel ist ein realistischer Überblick über Möglichkeiten, Risiken und Entwicklungsherausforderungen.}

\TopicLeitfragen{\item Welche Signalerfassungswege sind wofür geeignet?
\item Wie balanciert man Nutzen, Risiken und Ethik?
\item Welche Pfade in die Anwendung sind realistisch?
}
\TopicScope{\item Keine Medizinprodukte-Regulierung im Detail.
\item Keine Neurophysik im Beweisniveau.
}
\TopicPractice{Optionale Praxisidee: Architektur-Skizze eines BCI-Use-Case (Signalweg + Safety/Ethik-Check).}

\TopicKeywords{Neuroprothesen, BCI, Stimulation, Signalverarbeitung, Ethik}

\nocite{bciSurvey,whoRehab}
\end{topic}
